% TEMPLATE-MILC-v3.tex - plantilla maestra UNICA de las guias MILC v3.
%
% La usan las tres familias de guias, cada una con su driver:
%   · Grados 6-11  -> scripts/build-guias-g11.py
%   · Bebras       -> scripts/build-guias-bebras.py
%   · Semillero    -> scripts/build-guias-semillero.py
%
% Antes habia tres copias casi identicas (~400 lineas iguales, 6 distintas):
% cada mejora habia que aplicarla tres veces, y en la practica no se hacia
% (el motor de recursos vivia solo en dos de ellas). Lo que varia entre
% familias esta parametrizado en 6 placeholders que cada driver llena
% (se nombran aqui SIN su sintaxis real: el builder reemplaza por texto plano,
% tambien dentro de los comentarios, y un valor multilinea rompe el preambulo):
%
%   HEADER_IZQ        encabezado izquierdo de pagina
%   HEADER_DER        encabezado derecho de pagina
%   PORTADA_ETIQUETA  rotulo sobre el numero grande (GRADO/MOMENTO/MODULO)
%   PORTADA_SERIE     linea de serie de la portada
%   PORTADA_ID        identificador de la guia en la portada
%   PORTADA_CIRCULO   el circulito de grado (°), o vacio si no aplica
%
% El resto de claves las documenta content/guias/_SCHEMA.md. El script valida
% que no queden placeholders sin reemplazar antes de escribir el archivo final.

\documentclass[11pt,letterpaper]{article}
\usepackage[margin=2.0cm]{geometry}
\usepackage[table]{xcolor}
\usepackage{fontspec}
\usepackage[spanish]{babel}
\usepackage{tikz}
% Librerías TikZ para los diagramas de las guías (bloque `recursos:`):
%  · babel  → babel en español vuelve activos caracteres como «>», y TikZ los usa
%             en sus opciones (>=stealth, ->). Sin esta librería, cualquier
%             diagrama con flechas revienta con «\language@active@arg> has an
%             extra }». Debe ir ANTES de que se dibuje nada.
%  · positioning        → sintaxis «below left=2pt and 3pt».
%  · decorations.pathreplacing → llaves ({brace}) para señalar rangos.
\usetikzlibrary{babel,positioning,decorations.pathreplacing}
\usepackage{graphicx}
% Circuitos: el trabajo de electrónica se hace hoy con micro:bit y sus sensores
% integrados (MakeCode, sin cableado), así que no se necesitan esquemáticos.
% Si algún día entran montajes con protoboard, basta descomentar esta línea:
% los diagramas .tex/.tikz declarados en `recursos:` ya se incrustan solos.
% \usepackage{circuitikz}
\usepackage[most]{tcolorbox}
\usepackage{tabularx}
\usepackage{array}
\usepackage{fancyhdr}
\usepackage{titlesec}
\usepackage[document]{ragged2e}  % bandera a la derecha en todo el cuerpo
\usepackage{enumitem}
\usepackage{microtype}

% Helvetica Neue no trae la flecha U+2192, asi que cada «→» escrito literalmente
% en el YAML se compilaba a NADA, en silencio (462 flechas en 61 guias: rutas del
% tipo «paleta → tipografia → lockup» salian rotas). Mapeamos el caracter a su
% version matematica, que si tiene glifo. Asi el autor puede seguir escribiendo
% «→» comodamente en el YAML.
\usepackage{newunicodechar}
\newunicodechar{→}{\ensuremath{\rightarrow}}
% Mismo problema con los simbolos de marca y vinetas: el «✓» de «una palomita ✓»
% salia como cuadrito vacio en el PDF de todas las guias que lo usaban.
\usepackage{pifont}
\newunicodechar{✓}{\ding{51}}
\newunicodechar{✗}{\ding{55}}
\newunicodechar{✔}{\ding{52}}
\newunicodechar{•}{\textbullet}
\newunicodechar{≥}{\ensuremath{\geq}}
\newunicodechar{≤}{\ensuremath{\leq}}

\setmainfont{Helvetica Neue}
\setsansfont{Helvetica Neue}
\setlength{\parindent}{0pt}
\setlength{\parskip}{6pt}
% Sin particion de palabras: en 6.º-7.º una palabra cortada por guion al final
% de linea («Comuni-cacion») frena la lectura mucho mas que una linea corta.
\hyphenpenalty=10000
\exhyphenpenalty=10000
\pretolerance=2000
\tolerance=3000
\setlength{\emergencystretch}{3em}
\linespread{1.10}
\renewcommand{\arraystretch}{1.25}
\setlist{leftmargin=5mm,itemsep=1mm,topsep=1mm}
\newcolumntype{Y}{>{\RaggedRight\arraybackslash}X}

\definecolor{milcMagenta}{HTML}{E6007E}
\definecolor{milcVino}{HTML}{5A0038}
\definecolor{milcTurquesa}{HTML}{0093A5}
\definecolor{milcVerde}{HTML}{58A923}
\definecolor{milcMostaza}{HTML}{E5B400}
\definecolor{milcOcre}{HTML}{B86600}
\definecolor{milcNegro}{HTML}{241816}
\definecolor{milcGris}{HTML}{F7F7F4}
\definecolor{milcLinea}{HTML}{E8E2DD}
\definecolor{milcGradePrimary}{HTML}{84CC16}
\definecolor{milcGradeDark}{HTML}{4D7C0F}
\definecolor{milcGradeSoft}{HTML}{ECFCCB}
\definecolor{milcGradeAccent}{HTML}{CFE7FF}
\definecolor{milcGradeText}{HTML}{FFFFFF}
\color{milcNegro}

\pagestyle{fancy}
\fancyhf{}
\lhead{\small Tecnología e Informática · Grado 7}
\rhead{\small Guía 9 · MILC v3}
\cfoot{\small\thepage}
\renewcommand{\headrulewidth}{0pt}

\newtcolorbox{titlebox}[1]{
  enhanced,colback=#1,colframe=#1,arc=7mm,boxrule=0pt,
  left=7mm,right=7mm,top=3mm,bottom=3mm,
  before skip=8pt,after skip=8pt,
  fontupper=\sffamily\bfseries\Large\color{white}
}

\newtcolorbox{softbox}[3]{
  enhanced,colback=#2,colframe=#1,borderline west={4pt}{0pt}{#1},
  arc=2mm,boxrule=.4pt,left=4mm,right=4mm,top=3mm,bottom=3mm,
  before skip=6pt,after skip=8pt,title={#3},coltitle=white,
  colbacktitle=#1,fonttitle=\sffamily\bfseries,fontupper=\RaggedRight,
  attach boxed title to top left={xshift=3mm,yshift=-2mm},
  boxed title style={arc=2mm, boxrule=0pt}
}

\newtcolorbox{drawbox}[2]{
  enhanced,colback=white,colframe=#1,arc=4mm,boxrule=1pt,
  left=5mm,right=5mm,top=4mm,bottom=4mm,before skip=8pt,after skip=8pt,
  title={#2},coltitle=white,colbacktitle=#1,fonttitle=\sffamily\bfseries,
  fontupper=\RaggedRight,
  attach boxed title to top left={xshift=4mm,yshift=-2mm},
  boxed title style={arc=3mm, boxrule=0pt}
}

\newtcolorbox{infoband}[1]{
  enhanced,colback=#1!8,colframe=#1!50,arc=2mm,boxrule=.5pt,
  left=4mm,right=4mm,top=2mm,bottom=2mm,before skip=4pt,after skip=4pt,
  fontupper=\small\RaggedRight
}

% Puente narrativo entre secciones (contrato editorial Conéctate).
% Da continuidad: "Acabas de X. Ahora vas a Y porque Z."
% Visualmente: banda izquierda en color del grado, texto en itálica pequeña.
\newtcolorbox{puentebox}{
  enhanced,colback=milcGradeSoft,colframe=milcGradeDark,
  borderline west={3pt}{0pt}{milcGradeDark},
  arc=0mm,boxrule=0pt,left=5mm,right=4mm,top=2.5mm,bottom=2.5mm,
  before skip=4pt,after skip=8pt,
  fontupper=\itshape\small\color{milcNegro}\RaggedRight
}
% La flecha va en modo matematico a proposito: Helvetica Neue NO tiene el glifo
% U+2192, asi que \textrightarrow se compilaba a NADA («Missing character: There
% is no → in font Helvetica Neue») y el puente salia sin flecha. En modo
% matematico la flecha viene de la fuente matematica y si se dibuja.
\newcommand{\puente}[1]{%
  \begin{puentebox}%
    \textcolor{milcGradeDark}{$\rightarrow$}\hspace{2mm}#1%
  \end{puentebox}%
}

\newcommand{\linead}[1]{\noindent\color{milcLinea}\rule{#1}{0.5pt}\color{milcNegro}}
\newcommand{\respuesta}[1]{\par\vspace{2mm}\foreach \n in {1,...,#1}{\noindent\color{milcLinea}\rule{\linewidth}{0.5pt}\par\vspace{5mm}}\color{milcNegro}}
\newcommand{\checkbox}{\textcolor{milcGradeDark}{\fbox{\phantom{x}}}\hspace{5pt}}

\newcommand{\phasecard}[4]{
\begin{tcolorbox}[enhanced,colback=#3,colframe=#2,arc=3mm,boxrule=.5pt,height=3.05cm,valign=center,halign=center,left=1.5mm,right=1.5mm,top=2mm,bottom=2mm]
{\sffamily\bfseries\fontsize{22}{24}\selectfont\textcolor{#2}{#1}}\par
{\sffamily\bfseries\small #4}
\end{tcolorbox}
}

% ── Piezas del rediseno 2026 (legibilidad 6.º-11.º) ───────────────────
% Banda de estacion: reemplaza el titlebox macizo. Titulo a la izquierda,
% dato practico (tiempo/modalidad) a la derecha, en una sola linea.
\newcommand{\milcestacion}[2]{%
\begin{tcolorbox}[enhanced,colback=#1,colframe=#1,arc=2mm,boxrule=0pt,
  left=5mm,right=5mm,top=2.6mm,bottom=2.6mm,before skip=11pt,after skip=7pt]
{\sffamily\bfseries\fontsize{15}{18}\selectfont\color{white}#2}
\end{tcolorbox}}

% Tarjeta de paso numerada: sustituye las tablas «Pilar / Aplicacion», que
% eran formato de consulta docente, no de estudiante. El numero va en un
% circulo sobre el borde para que la secuencia se lea de un vistazo.
\newcommand{\milcpaso}[3]{%
\begin{tcolorbox}[enhanced,colback=white,colframe=milcLinea,arc=1.5mm,boxrule=.9pt,
  left=14mm,right=4mm,top=3mm,bottom=3mm,before skip=4pt,after skip=4pt,
  fontupper=\RaggedRight,
  overlay={\node[anchor=north west,circle,fill=#2,text=white,inner sep=0pt,
    minimum size=7.4mm,font=\sffamily\bfseries\small]
    at ([xshift=3.6mm,yshift=-3.2mm]frame.north west) {#1};}]
#3
\end{tcolorbox}}

% Casilla marcable de tamano util con lapiz (la anterior era \fbox{\phantom{x}},
% de alto variable segun el texto que la seguia).
\newcommand{\milcmarca}{\framebox[3.8mm]{\rule{0pt}{3.4mm}}}

% Fila marcable de lista suelta (checks de la estacion 1).
\newcommand{\milccheckrow}[1]{%
\par\vspace{1.8mm}\noindent
\makebox[6.4mm][l]{\raisebox{-.55mm}{\textcolor{milcNegro}{\milcmarca}}}%
\parbox[t]{\dimexpr\linewidth-6.4mm\relax}{\RaggedRight #1}\par}

% Celda marcable dentro de la rubrica (tabla de autoevaluacion). Misma
% sangria francesa, pero pensada para vivir dentro de una columna X.
\newcommand{\milcopcion}[1]{%
\makebox[5.2mm][l]{\raisebox{-.55mm}{\textcolor{milcNegro}{\milcmarca}}}%
\parbox[t]{\dimexpr\linewidth-5.2mm\relax}{\RaggedRight #1}}

% ── Recursos visuales de la guía ──────────────────────────────────────
% Declarados en el bloque `recursos:` del YAML y emitidos por el builder.
% \guiaFigura   → imagen rasterizada o PDF (foto, carta celeste, montaje).
% \guiaDiagrama → fuente TikZ/circuitikz (.tex) incrustada como vector:
%                 es la vía para circuitos y esquemáticos editables.
% #1 = ruta relativa al .tex · #2 = pie (puede ir vacío)
\newcommand{\guiaFigura}[2]{%
\begin{center}
\includegraphics[width=0.88\linewidth,height=0.38\textheight,keepaspectratio]{#1}\\[1.5mm]
{\footnotesize\itshape\textcolor{milcNegro!75}{#2}}
\end{center}
}

\newcommand{\guiaDiagrama}[2]{%
\begin{center}
\input{#1}\\[1.5mm]
{\footnotesize\itshape\textcolor{milcNegro!75}{#2}}
\end{center}
}

\begin{document}
\thispagestyle{empty}

% ── Portada ───────────────────────────────────────────────────────────
% Antes: un rectangulo de color a pagina completa. Se veia bien en pantalla,
% pero en la fotocopiadora del colegio se comia el toner y salia gris sucio.
% Ahora la banda ocupa el tercio superior y el resto va en blanco.
\begin{tikzpicture}[remember picture,overlay]
  \path[fill=milcGradePrimary]
    (current page.north west) rectangle ([yshift=-10.6cm]current page.north east);

  \node[anchor=north west,text=milcGradeText] at ([xshift=2.0cm,yshift=-1.55cm]current page.north west)
    {\sffamily\bfseries\fontsize{12}{14}\selectfont GRADO};
  \node[anchor=north west,text=milcGradeText,opacity=.85] at ([xshift=2.0cm,yshift=-2.35cm]current page.north west)
    {\sffamily\bfseries\fontsize{8.5}{10}\selectfont SERIE GUÍAS · TECNOLOGÍA E INFORMÁTICA};
  \node[anchor=north east,text=milcGradeText,opacity=.85,align=right] at ([xshift=-2.0cm,yshift=-1.55cm]current page.north east)
    {\sffamily\bfseries\fontsize{9}{12}\selectfont I.E. Sor María Juliana\\Cartago, Valle del Cauca};

  \node[anchor=west,text=milcGradeText] (gradonum) at ([xshift=1.85cm,yshift=-7.1cm]current page.north west)
    {\sffamily\bfseries\fontsize{112}{112}\selectfont 7};
  % El circulito de grado (°) solo aplica a las guias de grado («11°»); en
  % Bebras y Semillero el numero grande es un momento/modulo, no un grado.
  % Se posiciona relativo al nodo `gradonum`, no en coordenadas absolutas.
  \draw[milcGradeText,line width=1.6pt]
    ([xshift=2.0mm,yshift=-4.5mm]gradonum.north east) circle (.15cm);

  \node[anchor=west,text=milcGradeText,text width=11.4cm,align=left] at ([xshift=1.5cm,yshift=0.5cm]gradonum.east)
    {
     {\sffamily\bfseries\fontsize{9}{11}\selectfont\MakeUppercase{Período 1 · Trabajo colaborativo en la nube}}\\[3mm]
     {\sffamily\bfseries\fontsize{21}{25}\selectfont\setlength{\baselineskip}{27pt}Outlook y\\[4pt]Teams ---\\[4pt]comunicarme\\[4pt]con voz\\[4pt]profesional}\\[3.5mm]
     {\sffamily\bfseries\fontsize{15}{18}\selectfont Guía 9-7-TIC}
    };
\end{tikzpicture}

\vspace*{9.4cm}
\vfill

{\sffamily\bfseries\fontsize{8}{10}\selectfont\textcolor{milcGradeDark}{LO QUE TE LLEVAS HOY}}

\begin{tcolorbox}[enhanced,colback=white,colframe=milcNegro,arc=2mm,boxrule=1.6pt,
  left=5mm,right=5mm,top=4mm,bottom=4mm,before skip=3pt,after skip=12pt,
  fontupper=\RaggedRight\fontsize{11}{15.5}\selectfont]
\textbf{(a) 3 correos profesionales reales enviados desde Outlook} institucional (uno al profe, uno al coordinador, uno a un compañero); \textbf{(b) 1 reunión virtual programada} en Teams con tu equipo de la S6 (un equipo de 3); \textbf{(c) en el cuaderno}: tabla comparativa Outlook vs Teams vs WhatsApp con 5 criterios + bitácora de la reunión virtual.
\end{tcolorbox}

\begin{tcolorbox}[enhanced,colback=milcGris,colframe=milcGris,arc=2mm,boxrule=0pt,
  left=5mm,right=5mm,top=3.5mm,bottom=3.5mm,after skip=12pt]
{\sffamily\bfseries\fontsize{8}{10}\selectfont\textcolor{milcNegro!55}{ESTA GUÍA ES DE}}\\[3mm]
\begin{tabularx}{\linewidth}{X p{3.2cm} p{3.2cm}}
\linead{\linewidth} & \linead{3.0cm} & \linead{3.0cm}\\[-1mm]
{\fontsize{8.5}{10}\selectfont\textcolor{milcNegro!55}{Nombre completo}} &
{\fontsize{8.5}{10}\selectfont\textcolor{milcNegro!55}{Grupo}} &
{\fontsize{8.5}{10}\selectfont\textcolor{milcNegro!55}{Fecha}}\\
\end{tabularx}
\end{tcolorbox}

\vfill

\noindent\textcolor{milcLinea}{\rule{\linewidth}{1pt}}\\[2mm]
{\sffamily\bfseries\fontsize{9.5}{12}\selectfont Autor: Dr. Álvaro Cárdenas Orozco}\\[1mm]
{\sffamily\fontsize{8.5}{11}\selectfont\textcolor{milcNegro!55}{Metodología MILC · apertura ancestral · pensamiento computacional · triángulo Dussel–estoicismo–Floridi}}

\newpage

\section*{Apertura: Outlook y Teams --- comunicarme con voz profesional}

\begin{softbox}{milcGradeDark}{milcGradeSoft}{Saber ancestral}
\textbf{Doña Mercedes la maestra rural de Cartago tenía dos formas distintas de comunicar con la comunidad de la vereda La Plata.} Para asuntos \textbf{formales} (anuncios al rector, peticiones al alcalde, cartas de recomendación de sus alumnos), \emph{escribía una carta a mano} con letra clara, sobre, sello y firma. Tomaba 1 día en hacerla bien. Para asuntos \textbf{cotidianos} (avisar a una mamá que su hijo se enfermó, coordinar quién traía cuál material), iba a su casa o mandaba un \emph{recado verbal} con otro niño. \textbf{Eran 2 canales distintos para situaciones distintas}: la carta para lo serio que queda registrado; el recado para lo rápido que no necesita papel. \emph{Confundir los dos era ridículo}: nadie escribía una carta formal para decir \emph{``traiga lápiz mañana''}, ni mandaba recado verbal para una petición al alcalde. Hoy con tu cuenta institucional tienes los mismos 2 canales: \textbf{Outlook} (la carta digital) y \textbf{Teams} (el recado digital).
\end{softbox}

\begin{softbox}{milcTurquesa}{milcGris}{Saber contemporáneo}
\textbf{Outlook} es el cliente de correo electrónico institucional de Microsoft. Lo usas para \textbf{comunicación formal asincrónica}: la persona te responde cuando pueda, no de inmediato. Mensaje queda guardado, firmado, con tu nombre y rastro. Ideal para: \emph{mensajes a profesores y autoridades, solicitudes oficiales, peticiones a otros colegios o universidades, cualquier asunto serio que pueda necesitarse después como evidencia}. \textbf{Teams} es la app de comunicación profesional en tiempo real. Combina \emph{chat informal entre compañeros de equipo} (similar a WhatsApp pero institucional), \emph{llamadas y videollamadas}, \emph{reuniones programadas} con varios participantes, \emph{compartir pantalla} mientras hablas. Ideal para: \emph{coordinar trabajos en equipo, hacer reuniones virtuales, preguntar dudas rápidas, presentar avances al profe}. \textbf{Diferencia crucial:} \emph{Outlook = serio, formal, queda registrado}. \emph{Teams = ágil, conversacional, también queda pero menos formal}.
\end{softbox}

\begin{softbox}{milcMostaza}{milcGris}{Pregunta-puente}
\textit{Si tienes que \emph{pedirle al rector permiso para ir al baño durante un examen importante}, ¿\textbf{usas Outlook o Teams}? Y si tienes que \emph{preguntarle a tu compañero de equipo qué le pareció tu sección del trabajo grupal}, ¿usas Outlook o Teams?}
\end{softbox}

\begin{softbox}{milcVerde}{milcGris}{Saber y saber hacer}
Al terminar la clase: (1) podrás \textbf{identificar} la diferencia entre Outlook y Teams; (2) sabrás \textbf{aplicar} las 5 partes de un correo formal (ya aprendidas en G6); (3) podrás \textbf{programar} una reunión virtual en Teams; (4) habrás \textbf{enviado} 3 correos profesionales reales + 1 reunión Teams.
\end{softbox}



\small
\begin{tabularx}{\textwidth}{>{\bfseries}p{3.0cm}Y}
\rowcolor{milcGradePrimary!12}Autor & Dr. Álvaro Cárdenas Orozco\\
DBA & Utiliza herramientas digitales para trabajar de manera colaborativa con otros (MEN, T\&I 7°)\\
Referentes & Saber ancestral de la minga · Microsoft 365 y OneDrive · Coautoría asincrónica\\
Producto final & \textbf{(a) 3 correos profesionales reales enviados desde Outlook} institucional (uno al profe, uno al coordinador, uno a un compañero); \textbf{(b) 1 reunión virtual programada} en Teams con tu equipo de la S6 (un equipo de 3); \textbf{(c) en el cuaderno}: tabla comparativa Outlook vs Teams vs WhatsApp con 5 criterios + bitácora de la reunión virtual.\\
\end{tabularx}
\normalsize

\puente{Hoy haces 4 cosas: diferencias Outlook y Teams, envías 3 correos formales, programas una reunión Teams, evalúas comunicaciones reales.}

\milcestacion{milcGradeDark}{Ruta MILC: así vamos a aprender}
\begin{tabularx}{\textwidth}{>{\centering\arraybackslash}X>{\centering\arraybackslash}X>{\centering\arraybackslash}X>{\centering\arraybackslash}X}
\phasecard{1}{milcVerde}{milcGris}{Escucha\\[-1mm]\small Observo, escucho y pregunto.} &
\phasecard{2}{milcTurquesa}{milcGris}{Entender\\[-1mm]\small Sistematizo y ordeno.} &
\phasecard{3}{milcMagenta}{milcGris}{Hacer\\[-1mm]\small Construyo y aplico.} &
\phasecard{4}{milcVino}{milcGris}{Cierre\\[-1mm]\small Evalúo y reflexiono.}
\end{tabularx}


\puente{Empezamos con un test de 6 mensajes: ¿Outlook, Teams o WhatsApp?}

\milcestacion{milcVerde}{Estación 1 de 3 · Escucha}

\begin{softbox}{milcVerde}{milcGris}{Escena para escuchar}
\textbf{Actividad 1 · IDENTIFICA --- 6 mensajes, 3 canales} (10 min · individual). Lee estos 6 mensajes y escribe en el cuaderno qué canal usarías: \textbf{O (Outlook formal)}, \textbf{T (Teams)}, \textbf{W (WhatsApp personal)}. \emph{(1) Pedirle al rector que considere reabrir el cupo de un compañero. (2) Preguntarle a un amigo cercano si va a la fiesta del sábado. (3) Coordinar con tu equipo de 3 dónde se reúnen mañana para hacer el trabajo. (4) Solicitar a una universidad información sobre carreras. (5) Mandar un meme gracioso a tu mejor amiga. (6) Avisar a tu profe que vas a faltar mañana por cita médica.}
\end{softbox}

{\sffamily\bfseries\fontsize{8}{10}\selectfont\textcolor{milcNegro!55}{MARCA CUANDO LO HAYAS HECHO}}
\milccheckrow{Leí los 6 mensajes con calma.}
\milccheckrow{Marqué O, T o W en cada uno.}
\milccheckrow{Pensé en el nivel de formalidad de cada situación.}

\begin{infoband}{milcVerde}


\textbf{Lo que va al cuaderno (Actividad 1 · IDENTIFICA).} Título: «Actividad 1 --- 6 mensajes, 3 canales». \textbf{Formato:} lista numerada con O/T/W. \textbf{Extensión:} 6 líneas. \textbf{Sabes que terminaste cuando} cada mensaje tiene su canal asignado con criterio.
\end{infoband}

\puente{Después aprendes las 3 herramientas con sus usos correctos.}

\milcestacion{milcTurquesa}{Estación 2 de 3 · Entender}

\begin{softbox}{milcTurquesa}{milcGris}{Lo que vas a entender}
Tu intuición se afila. Aprende los \textbf{usos correctos de cada canal} y por qué Outlook + Teams son superiores a WhatsApp para asuntos académicos y profesionales. \emph{Respuestas correctas:} 1=O, 2=W, 3=T, 4=O, 5=W, 6=O.
\end{softbox}

\milcpaso{1}{milcTurquesa}{\textbf{Outlook --- el canal formal.} Lo usas para mensajes que \textbf{necesitan registro, autoridad y formalidad}. Tu cuenta institucional ya tiene Outlook activo en \texttt{outlook.office.com}. \textbf{Las 5 partes obligatorias del correo formal} (ya las viste en G6·P1·S3): \emph{(1) Destinatario}, \emph{(2) Asunto claro}, \emph{(3) Saludo formal}, \emph{(4) Cuerpo en 3 bloques}, \emph{(5) Despedida + firma con tu nombre + grado}. \textbf{Funciones útiles de Outlook:} \emph{calendario integrado} (programa reuniones desde el correo), \emph{firma automática} (al final de todos tus correos pones tu firma fija), \emph{categorías y reglas} (organiza correos por color). \textbf{Una excelente práctica:} crea una firma fija en Outlook con tu nombre + grado + colegio.}
\milcpaso{2}{milcTurquesa}{\textbf{Teams --- el canal ágil profesional.} Lo usas para \textbf{comunicación rápida entre miembros de un mismo equipo o curso}, sin la formalidad del correo. \textbf{3 componentes principales:} \emph{(a) Chat:} conversaciones por escrito (similar a WhatsApp pero institucional, queda guardado en la cuenta del colegio). \emph{(b) Llamadas y videollamadas:} 1 a 1 o en grupo, audio o video, con compartir pantalla. \emph{(c) Reuniones programadas:} eventos con hora, fecha, lista de invitados y enlace que todos abren al mismo tiempo. \textbf{Cómo programar reunión en Teams:} ícono \emph{Calendario} → \emph{Nueva reunión} → poner título, fecha, hora, agregar invitados → \emph{Enviar}. Los invitados reciben correo en Outlook con el enlace y se les agrega a su calendario automáticamente.}
\milcpaso{3}{milcTurquesa}{\textbf{WhatsApp vs Outlook vs Teams --- la tabla maestra.} \textbf{WhatsApp:} para amigos y familia. Personal. No queda en archivo institucional. NUNCA para asuntos académicos serios. \textbf{Outlook:} para asuntos formales. Asincrónico. Queda registrado. NO para coordinación rápida de equipos. \textbf{Teams:} para equipos profesionales. Sincrónico/asincrónico. Queda registrado en cuenta institucional. NO para asuntos personales con familia.}
\milcpaso{4}{milcTurquesa}{\textbf{Errores frecuentes al elegir canal.} (1) \emph{Mandar al profe un meme por WhatsApp en lugar de un correo formal por situación seria} --- mezclas niveles, generas confusión. (2) \emph{Hacer una solicitud al coordinador por chat de Teams en lugar de un correo Outlook} --- los chats se pierden rápido; el correo queda registrado. (3) \emph{Programar reunión Teams sin agenda} --- los participantes no saben qué se va a discutir; pierde tiempo. (4) \emph{No mirar el correo institucional regularmente} --- te pierdes anuncios importantes. Revisa mínimo 1 vez al día.}

\begin{drawbox}{milcTurquesa}{3 canales con sus usos}
\textbf{Outlook:} formal, registrado. Correos a profes, autoridades, instituciones. \\[1mm]
\textbf{Teams:} ágil profesional. Chats de equipo, reuniones, videollamadas. \\[1mm]
\textbf{WhatsApp:} personal. Familia, amigos. Nunca para asuntos académicos serios.
\end{drawbox}

\begin{softbox}{milcMostaza}{milcGris}{Errores comunes y su corrección}
\textbf{Mitos sobre Outlook y Teams.} (1) \emph{``Outlook es solo correo''} --- Falso. Trae calendario, contactos, tareas, reglas. (2) \emph{``Teams es lo mismo que WhatsApp''} --- Falso. Teams es institucional, queda en archivo del colegio, tiene reuniones programadas con calendario. (3) \emph{``Tengo que tener Outlook abierto siempre''} --- Falso. Lo abres 1-2 veces al día. Te llegan notificaciones cuando algo importante entra. (4) \emph{``Es mejor mandar todo por WhatsApp porque es más rápido''} --- Falso a largo plazo: los mensajes serios sin registro se pierden y no tienen peso.
\end{softbox}

\begin{infoband}{milcTurquesa}


\textbf{Lo que va al cuaderno (Actividad 2 · EVALÚA).} Título: «Actividad 2 --- Cuál canal corresponde a cuál situación». \textbf{Formato:} tabla de 5 filas, 3 columnas. Filas: \emph{(1) ¿Para qué?}, \emph{(2) Formalidad}, \emph{(3) Queda registrado}, \emph{(4) Ejemplo de uso}, \emph{(5) Cuándo NO usar}. \textbf{Veredicto (5 casos):} para estas 5 situaciones, decide qué canal corresponde y justifica con 1 razón: (1) pedir prórroga de un trabajo al profe; (2) coordinar quien lleva qué al asado del salón; (3) preguntar la temática del próximo parcial; (4) avisar que vas a llegar tarde a clase; (5) enviar tu hoja de vida para una práctica. \textbf{Extensión:} tabla + 5 veredictos. \textbf{Sabes que terminaste cuando} distingues canal formal (Outlook), canal de equipo (Teams) y canal informal (WhatsApp) sin dudar.
\end{infoband}

\puente{Con todo claro, envías los 3 correos y programas la reunión Teams.}

\milcestacion{milcMagenta}{Estación 3 de 3 · Hacer}

\begin{softbox}{milcMagenta}{milcGris}{Cómo vas a construir tu producto}
\textbf{Actividad 3 · APLICA --- 3 correos profesionales + 1 reunión Teams} (32 min · individual). Vas a enviar 3 correos reales por Outlook y programar 1 reunión de Teams con tu equipo de la S6. Si no tienes computador en clase, planeas todo en el cuaderno y ejecutas en casa esta semana.
\end{softbox}

\milcpaso{1}{milcMagenta}{\textbf{Paso 1 --- Abre Outlook.} Ve a \texttt{outlook.office.com}, inicia sesión con tu cuenta institucional. Verás la bandeja de entrada. Antes de empezar, ve a \emph{Configuración → Firma de correo electrónico} y crea tu firma fija: \emph{``Atentamente, [Tu nombre completo], Grado 7[letra], I.E. Sor María Juliana, año 2026.''}. Guarda. De ahora en adelante esa firma aparece automáticamente al final de cada correo que escribas.}
\milcpaso{2}{milcMagenta}{\textbf{Paso 2 --- Correo 1: al profe.} \emph{Nuevo correo}. Destinatario: el correo institucional de tu profe de Tecnología. Asunto: \emph{``Solicitud de aclaración sobre la actividad de la próxima sesión''}. Cuerpo en 3 bloques: \emph{(1) qué quieres aclarar (1 frase)}, \emph{(2) por qué tienes esa duda (1 frase)}, \emph{(3) qué pides (que te confirme algo, te explique, etc., 1 frase)}. Máximo 5 líneas en total. Tu firma aparece sola al final. Envía.}
\milcpaso{3}{milcMagenta}{\textbf{Paso 3 --- Correo 2: al coordinador.} Destinatario: correo del coordinador (pregunta al profe si no lo tienes). Asunto: \emph{``Sugerencia para la sala de Tecnología --- estudiante de grado 7''}. Cuerpo: una sugerencia constructiva (un ventilador, mejor luz, una pizarra adicional, lo que tú propongas). 3 bloques: idea, por qué la propones, qué necesita para implementarse. Envía.}
\milcpaso{4}{milcMagenta}{\textbf{Paso 4 --- Correo 3: a un compañero.} Destinatario: correo institucional de un compañero. Asunto: \emph{``Coordinación trabajo grupal G7-P1''}. Cuerpo: propón día y hora para que se reúnan a revisar el trabajo de la S6, o pregúntale algo concreto sobre el trabajo. Envía. Después verifica con él que recibió el correo.}
\milcpaso{5}{milcMagenta}{\textbf{Paso 5 --- Reunión Teams.} Ve al ícono de \emph{Teams} (calendario en la barra lateral). Clic en \emph{Nueva reunión}. Llena los datos: \emph{Título}: \emph{``Revisión trabajo grupal G7-P1''}. \emph{Fecha y hora}: alguna esta semana, 30 minutos. \emph{Agregar participantes}: los 2 compañeros de tu equipo de S6 (escribe sus correos). \emph{Agenda}: \emph{``(1) revisar trabajo colaborativo, (2) acordar cambios finales, (3) entregar al profe''}. Clic en \emph{Enviar}. Tus 2 compañeros recibirán invitación.}
\milcpaso{6}{milcMagenta}{\textbf{Paso 6 --- Bitácora + reflexión + firma.} En el cuaderno: \emph{(a) Lista de los 3 correos enviados (destinatario + asunto). (b) Datos de la reunión Teams (fecha, hora, participantes, agenda). (c) Reflexión: ``Mi descubrimiento de hoy fue: \_\_\_\_''}. Firma con nombre y fecha. Esta sesión te entrenó en comunicación profesional digital.}

\begin{drawbox}{milcMagenta}{Antes de cerrar la S9}
\checkbox Mi firma fija está configurada en Outlook. \\[1mm]
\checkbox Envié 3 correos profesionales (profe, coordinador, compañero). \\[1mm]
\checkbox Cada correo tiene las 5 partes y máximo 5 líneas de cuerpo. \\[1mm]
\checkbox Programé 1 reunión en Teams con mi equipo de 3. \\[1mm]
\checkbox La reunión tiene agenda con 3 puntos claros. \\[1mm]
\checkbox La bitácora está firmada con nombre y fecha.
\end{drawbox}

\begin{softbox}{milcMagenta}{milcGris}{Plantilla de trabajo}
\textbf{Estructura.} \textit{Paso 1:} firma fija. \textit{Paso 2:} correo 1 al profe. \textit{Paso 3:} correo 2 al coordinador. \textit{Paso 4:} correo 3 al compañero. \textit{Paso 5:} reunión Teams. \textit{Paso 6:} bitácora + firma.
\end{softbox}

\begin{infoband}{milcMagenta}


\textbf{Lo que va al cuaderno (Actividad 3 · APLICA).} Título: «Actividad 3 --- Mis primeros correos profesionales + reunión Teams». \textbf{Formato:} datos de los 3 correos + reunión Teams + reflexión. \textbf{Extensión:} 1 página. \textbf{Sabes que terminaste cuando} los 3 correos están enviados Y la reunión está programada con invitaciones recibidas por los compañeros.
\end{infoband}

\puente{El producto son los 3 correos enviados + reunión Teams + tabla comparativa.}

\milcestacion{milcMostaza}{Producto de la sesión}
\textbf{3 correos profesionales enviados + 1 reunión Teams programada · bitácora firmada}.
\begin{softbox}{milcMostaza}{milcGris}{Criterios de calidad}
(1) Los 3 correos siguen las 5 partes obligatorias. (2) Cada correo tiene cuerpo máximo de 5 líneas. (3) La firma fija aparece al final de cada correo. (4) La reunión Teams tiene título, fecha, hora, 2 invitados, agenda. (5) La bitácora está completa y firmada.
\end{softbox}

\puente{Antes de salir, verifica que los 3 correos cumplen las 5 partes y que la reunión está programada con invitaciones enviadas.}

\milcestacion{milcVino}{Evaluación liberadora · cinco dimensiones}

\begin{softbox}{milcVino}{milcGris}{Sentido de la evaluación}
Evaluar no es solo revisar una nota. Es reconocer qué aprendí, qué cambió en mí, qué emociones manejé, cómo me relaciono con la ciudad y la comunidad, y qué le dejaré a quien viene detrás.
\end{softbox}

\textbf{Mi reflexión liberadora en cinco dimensiones.} Completa cada frase iniciada:

\small
\begin{tabularx}{\textwidth}{>{\bfseries\RaggedRight\arraybackslash}p{3.4cm}Y>{\RaggedRight\arraybackslash}p{4.6cm}}
\rowcolor{milcVino}\color{white}Dimensión & \color{white}\bfseries Mi respuesta (frame starter) & \color{white}\bfseries Algo que descubrí\\
1. Personal & Lo que más cambió en mí fue \linead{6.5cm} & \linead{4.2cm}\\
2. Emocional & La emoción más fuerte fue \linead{2.5cm} y la regulé \linead{3cm} & \linead{4.2cm}\\
3. Ciudadana & En democracia esto importa porque \linead{6cm} & \linead{4.2cm}\\
4. Local (Cartago) & En mi barrio sería útil para \linead{6.5cm} & \linead{4.2cm}\\
5. Intergeneracional & A mi abuela/abuelo le diría \linead{6.5cm} & \linead{4.2cm}\\
\end{tabularx}
\normalsize

\begin{infoband}{milcVino}
\textbf{Para la próxima sustentación.} Vuelve a esta tabla en seis meses. Verás que las respuestas cambian — no porque lo aprendido cambie, sino porque tú cambias. La evaluación liberadora es una conversación contigo a través del tiempo.
\end{infoband}


\puente{Cerramos con 3 voces sobre por qué usar el canal correcto es respetar al lector.}

\milcestacion{milcGradeDark}{Triángulo de pensamiento · cierre filosófico}

\begin{softbox}{milcVino}{milcGris}{Dussel · lente del nosotros}
\textit{Hablar igual con todos no es respeto: es desconocimiento. El respeto adapta el lenguaje al lector.}

{\footnotesize\itshape\color{milcNegro!70}--- formulación de esta guía, al modo de Enrique Dussel. No es una cita textual.}

Hay gente que cree que \emph{``yo hablo igual con todo el mundo''} es virtud. No es: es \textbf{confusión de canales}. \emph{Hablarle al rector como le hablas a tu amigo es no entender que el rector ocupa un lugar distinto en la conversación}. Doña Mercedes mandaba carta formal al alcalde y recado verbal al vecino: \emph{2 canales, mismas reglas de respeto adaptadas al contexto}. Hoy en digital, Outlook es para uno; Teams para otro; WhatsApp para otro más. Saber elegir es respeto.

\textbf{¿Estoy hablando igual con todos sin distinguir contextos? ¿Eso me ayuda o me complica?}
\end{softbox}
\linead{\linewidth}\\[1mm]
\linead{\linewidth}

\begin{softbox}{milcMostaza}{milcGris}{Estoicismo · Marco Aurelio · cuidado interior}
\textit{La carta formal queda. El recado se va con el viento. Sabe usar las dos a su tiempo.}

{\footnotesize\itshape\color{milcNegro!70}--- formulación de esta guía, al modo de Marco Aurelio. No es una cita textual.}

Marco Aurelio gobernaba el imperio con \textbf{dos canales}: \emph{decretos escritos} (que quedaban en archivo para siempre) y \emph{conversaciones verbales} (rápidas, ágiles, no registradas). Sabía cuándo usar cada uno. En el digital de hoy, esos dos canales son Outlook y Teams. Quien conoce ambos opera con \emph{precisión profesional}; quien solo conoce uno (o usa WhatsApp para todo) opera con menos finura.

\textbf{¿Estoy desarrollando precisión profesional en comunicación, o me quedo en informalidad permanente?}
\end{softbox}
\linead{\linewidth}\\[1mm]
\linead{\linewidth}

\begin{softbox}{milcTurquesa}{milcGris}{Floridi · lente de la infoesfera}
\textit{Tener correo institucional y saber usarlo es el primer rito de paso a la vida adulta digital.}

{\footnotesize\itshape\color{milcNegro!70}--- formulación de esta guía, al modo de Luciano Floridi. No es una cita textual.}

\textbf{Tu correo institucional} (\texttt{...@sormariajuliana.edu.co}) es \emph{tu primera identidad oficial digital}. Lo vas a reemplazar después por \emph{...@universidad.edu.co} y luego por \emph{...@empresa.com}. Pero el oficio es el mismo: \emph{escribir formal, claro, con firma, con asunto, con respeto}. Hoy entrenas ese oficio en la versión escolar. Mañana lo usarás en la versión universitaria. En 10 años en la profesional. Lo que aprendes hoy queda.

\textbf{¿Estoy aprovechando mi correo institucional como entrenamiento para mi yo de 25 años?}
\end{softbox}
\linead{\linewidth}\\[1mm]
\linead{\linewidth}

\begin{infoband}{milcNegro}
\textbf{Nota para el docente.} Las tres formulaciones son del autor de la guía, no citas textuales de los pensadores: recogen su lente para pensar el tema, no sus palabras. Si vas a citarlos en clase, remite a la obra original.
\end{infoband}

\milcestacion{milcVino}{Mi autoevaluación · marca una casilla por fila}

{\small
\setlength{\extrarowheight}{2.2pt}
\begin{tabularx}{\textwidth}{>{\RaggedRight\arraybackslash\bfseries}p{3.0cm}YYY}
\rowcolor{milcVino}\color{white}\bfseries Criterio & \color{white}\bfseries Lo logré & \color{white}\bfseries En proceso & \color{white}\bfseries Necesito apoyo\\[.6mm]
Comprensión del tema & \milcopcion{Explico las ideas con mis palabras.} & \milcopcion{Reconozco las ideas, falta precisión.} & \milcopcion{Necesito repasar lo básico.}\\[.8mm]
\rowcolor{milcGris}Pensamiento computacional & \milcopcion{Usé los pilares al armar mi producto.} & \milcopcion{Usé algunos pilares.} & \milcopcion{Necesito releer la estación 2.}\\[.8mm]
Aplicación y producto & \milcopcion{Mi producto cumple los criterios.} & \milcopcion{Está armado, con ajustes pendientes.} & \milcopcion{Aún está en construcción.}\\[.8mm]
\rowcolor{milcGris}Reflexión 5 dimensiones & \milcopcion{Respondí cada una con honestidad.} & \milcopcion{Respondí algunas con detalle.} & \milcopcion{Mis respuestas son muy generales.}\\[.8mm]
Triángulo de pensamiento & \milcopcion{Conecté las 3 voces con mi trabajo.} & \milcopcion{Conecté al menos una.} & \milcopcion{No las relaciono con mi proceso.}\\[.8mm]
\end{tabularx}}

\vspace{3mm}
\textbf{Hoy entendí que:}\\[1mm]
\linead{\linewidth}\\[1mm]
\linead{\linewidth}

\textbf{Mi compromiso para la próxima sesión es:}\\[1mm]
\linead{\linewidth}\\[1mm]
\linead{\linewidth}

\begin{softbox}{milcMostaza}{milcGris}{Cita de despedida · Marco Aurelio}
\textit{``El obstáculo en el camino se vuelve el camino. Lo que estorba al camino se vuelve camino.''}\\[1mm]
Cada error de hoy, bien observado, se convierte en pieza del camino que pisarás mañana.
\end{softbox}

\small
\begin{tabularx}{\textwidth}{>{\bfseries}p{3.6cm}Y}
\rowcolor{milcGradeDark!12}Nombre completo: & \linead{10cm}\\
Fecha: & \linead{4cm} \quad Curso/grupo: \linead{4cm}\\
Mi firma: & \linead{9cm}\\
\end{tabularx}
\normalsize

\begin{infoband}{milcGradeDark}
\textbf{Para la próxima sesión.} Esta semana, revisa tu Outlook al menos cada 2 días. Responde lo que tengas pendiente. Asiste a la reunión Teams que programaste. La próxima sesión es \textbf{Sesión 10 --- la cosecha del periodo}: tu equipo de 3 entrega un mini-proyecto colaborativo que reúne TODO lo aprendido en P1. Va a ser el cierre del periodo.
\end{infoband}

\end{document}
